\subsection{Filtered totalizations of Kan simplex categories}

\begin{theorem}[Proved: contractible filtered Kan totalizations are exact]
\label{C:kan-filtered-totalization}
Let $I$ be a nonempty small filtered category, let $K_i$ be nonempty
connected small Kan complexes, and let
\[
 A_i=\Delta_{\mathrm{inj}}/K_i\quad(i\in I)
\]
form a diagram of categories with arbitrary transition functors.
The transition functors need not be faithful or induced by simplicial
maps. Put $A=\colim_{i\in I}A_i$, the colimit in small categories.
Then the following conditions are equivalent:
\[
 BA\simeq *;\qquad A\text{ is total};\qquad
 \colim_A\text{ is exact over }\Z.
\]
When these conditions hold, $\colim_A$ is exact over every unital ring.
Consequently no total counterexample to C.4, and no total nonexact
category, can be obtained as such a filtered colimit. This includes
arbitrary coefficient groups and arbitrary connected Kan stages;
it imposes no increasing-connectivity condition on the stages.
\end{theorem}
\begin{proof}
We prove the nontrivial implication from weak contractibility of $BA$
to integral exactness. We first give the finite obstruction and its
functorial comparison, including the distinction between ordinary
and fat realization.

\emph{A finite family and a contractible path-space lift.}
Write $A(K)=\Delta_{\mathrm{inj}}/K$. For a finite parallel family
$f=(a_0,\ldots,a_m):P\rightrightarrows Q$ in $A(K)$, put
\[
 R_f=\operatorname{coeq}\left(
 \coprod_{j=1}^{m}\Delta^{P}\;
 \substack{\longrightarrow\\[-5pt]\longrightarrow}\;
 \Delta^{Q}\right).
 \tag{KC.1}
\]
Here $\Delta^B$ is the standard simplicial simplex on the finite
ordered vertex set of $B$, and the two arrows on the $j$th summand
are $a_j$ and $a_0$. The coequalizer is taken in simplicial sets.
It has finitely many nondegenerate simplices, although vertices or
faces can be identified. Its realization is therefore a finite CW
complex. All generating identifications are induced by injective
face maps and preserve degeneracy type, so a nondegenerate cell is
not identified with a degenerate simplex. The simplex label of $Q$ induces a simplicial map
$c_f:R_f\to K$.

If $|c_f|$ is nullhomotopic, the family $f$ has an integral affine
equalizer. Indeed, choose a vertex $k\in K$ and the based path-space
Kan fibration $\pi:P_kK\to K$. Its total space is weakly contractible;
the precise construction and proof are
\cite[Lemma I.7.5, pp.~30--31]{ABGoerssJardine}. The map
\[
 \operatorname{Map}(R_f,P_kK)\longrightarrow
 \operatorname{Map}(R_f,K)
\]
is a Kan fibration by \cite[Proposition I.5.2 and Corollary I.5.3]{ABGoerssJardine}.
Its image on vertices is a union of path components: lift an edge
starting at a lifted vertex, and repeat along an edge path. The
constant map at $k$ lifts, and $c_f$ belongs to its component.
Thus $c_f$ has a strict simplicial lift to $P_kK$. Pulling this lift
back to $\Delta^P$ and $\Delta^Q$ gives objects $P^*,Q^*$ and the
entire parallel family in $A(P_kK)$. This category is Ramsey and
integrally exact by Theorem~\ref{C4:contractible-kan-ramsey}.
An integral affine equalizer there is carried to one in $A(K)$ by
$\pi$. All augmentations are preserved. We need only this sufficient
nullhomotopy criterion; we do not assert its converse for arbitrary
$K$.

\emph{The map of base spaces associated with an arbitrary functor.}
Let $F:A(K)\to A(L)$, where $K$ is simply connected and $L$ is
connected. Denote by $\|K\|$ the fat realization retaining degenerate
simplices as cells. Then
\[
 \|K\|=\colim_{B\in A(K)}|\Delta^B|,
 \qquad BA(K)\simeq\|K\|\simeq |K|.
 \tag{KC.2}
\]
For clarity, the first weak equivalence can be obtained from the
homotopy colimit of the displayed simplex diagram. Its augmentation
to $BA(K)$ is a weak equivalence: in each bar degree it is the
disjoint union, over strings of arrows, of the weak equivalences
from a simplex to a point. Its map to the
ordinary colimit is a weak equivalence: order the index by simplex
dimension; at each object the colimit of its proper faces is exactly
the boundary of that simplex, so the diagram is built by cofibrations.
The second equivalence is the degeneracy comparison already checked
in Theorem~\ref{B:kan-recovery}. This also specifies the comparison
used below; no invariance under an arbitrary functor is presumed.

For every vertex label $v\in K_0$, choose a vertex $z_v$ of the
nonempty object $F(\{v\})$. For every object $B$, the assignments
\[
 j_B(x)=F(\{x\}\hookrightarrow B)(z_{\operatorname{label}(x)})
 \quad(x\in B)
\]
extend affinely to continuous maps
$j_B:|\Delta^B|\to|\Delta^{F B}|$. They are natural for all arrows
of $A(K)$. They need not preserve vertex orders, which causes no
problem for maps of geometric simplices. Composing with the target
simplex labels gives
\[
 J_F:\|K\|\longrightarrow\|L\|.
 \tag{KC.3}
\]
The same natural maps on homotopy colimits, followed by the change of
index along $F$, lie over $BF$ on the base augmentations. Hence
$J_F$ represents $BF$ under (KC.2). The canonical map
$\bar c_f:|R_f|\to\|K\|$ is obtained from the cell of $Q$;
its composite to $|K|$ is $|c_f|$.

We claim that if $J_F\bar c_f$ is nullhomotopic, then $Ff$ has an
integral affine equalizer. First suppose $K$ is countable and use
its homogeneous limit $M$ and automorphism group $G$. Theorem
\ref{B:kan-recovery} gives strong amalgamation, (St), (Op), and
$\operatorname{Fix}_M(G_B)=B$ for every finite $B\subset M$.
Every finite subset is an admissible substructure. Form the smooth
ordered $G$-set
\[
 S=\colim_{B\subset M,\ B\ne\varnothing}
       V(FB),
\]
where this colimit uses only inclusions. The maps are injective, so
we regard the vertex sets as subsets of $S$. Every finite subset of
$S$ lies in some $V(FB)$. Their coherent labels give a continuous
$G$-invariant map $c_S:|\Delta^S|\to\|L\|$.

If $S$ has a $G$-fixed vertex $z$, cone each point of
$|\Delta^{FQ}|$ linearly to $z$ in $|\Delta^S|$, then apply $c_S$.
This homotopy descends to $|R_{Ff}|$. Indeed, extend each
$a_j a_0^{-1}$ to an element $g_j\in G$; it sends the image of
$F a_0$ to that of $F a_j$, fixes $z$, and leaves $c_S$ invariant.
The resulting homotopy is a nullhomotopy of $\bar c_{Ff}$, so the
path-space criterion applies.

Suppose now that $S$ has no fixed vertex. By (Op), every $s\in S$
has a unique nonempty finite support $B_s$ with
$G_s=G_{B_s}$. If $s\in V(FB)$ then $G_B\leq G_s$, and
$\operatorname{Fix}_M(G_B)=B$ implies $B_s\subset B$.
The natural maps $j_B$ above assemble to a $G$-equivariant vertex
map $j:M\to S$. A point $j(x)$ has support exactly $\{x\}$:
its support is contained in that singleton and cannot be empty.
In particular $j$ is injective. Define the positive barycentric map
\[
 p:|\Delta^S|\longrightarrow|\Delta^M|,
 \qquad p(s)=\frac1{|B_s|}\sum_{x\in B_s}x.
 \tag{KC.4}
\]
This is a continuous affine map on every finite simplex, hence
continuous for the CW topology. It is $G$-equivariant and restricts
to $p_B:|\Delta^{FB}|\to|\Delta^B|$. These restrictions are natural
under every embedding, by extension to $G$. We have $p_Bj_B=1$,
and the straight-line homotopy
\[
 H_B(u,t)=(1-t)u+t\,j_Bp_B(u)
 \tag{KC.5}
\]
stays in $|\Delta^{FB}|$ and is natural under all arrows.
Consequently $p_B,j_B,H_B$ descend for the family to maps
$p_f:|R_{Ff}|\to|R_f|$, $j_f:|R_f|\to|R_{Ff}|$ and a homotopy
$j_fp_f\simeq1$. The label maps satisfy
\[
 \bar c_{Ff}\simeq J_F\bar c_f p_f.
 \tag{KC.6}
\]
Therefore the assumed nullhomotopy of $J_F\bar c_f$ gives a
nullhomotopy of $c_{Ff}$ after passage to ordinary realization,
and hence an integral affine equalizer.

The countability restriction on $K$ is unnecessary for this local
claim. Any specified finite family lies in $A(K')$ for a countable
simply connected Kan subcomplex $K'\subset K$: start with its
finitely many labels, close under degeneracies and faces, fill the
countably many horns, add paths and nullhomotopies for all edge
loops, and repeat. The ambient connected simply connected $K$
supplies these choices, and their countable union is Kan and simply
connected. Apply the preceding argument to $F|_{A(K')}$; its image
vertex set is countable regardless of the size of $L$. The finite
family, the chosen $j_B$, and its nullhomotopy in $\|L\|$ are
unchanged. This proves the claim for all small simply connected $K$.

\emph{Passage to a filtered stage.}
For a filtered diagram of small categories, the natural map
$\colim_i N A_i\to N(\colim_i A_i)$ is an isomorphism: each
composable finite string and each equality between two such strings
is represented at one stage. Homotopy classes from a finite
simplicial set commute with these filtered colimits after fibrant
replacement. One explicit verification uses $\operatorname{Ex}^\infty$:
it is a natural Kan replacement by
\cite[Theorem III.4.8, p.~188]{ABGoerssJardine}, and it commutes
with filtered colimits because each $\operatorname{Sd}\Delta^n$
is finite. A map from a finite simplicial set, and a homotopy from
it, therefore use only finitely many simplices and equalities and
are realized at one stage. This applies also to the finite sets
$R_f$ and to a simplicial circle.

Assume $BA\simeq *$ and take a finite parallel family in $A$.
It lifts to a parallel family $f:P\rightrightarrows Q$ in some
$A_i$. The finitely many loops $a_j a_0^{-1}$ in the fundamental
groupoid of $BA_i$ become trivial at one later stage $A_j$.
Replace $f$ by its image there. The universal simplicial cover
$\widetilde K_j\to K_j$ induces the universal covering category
$A(\widetilde K_j)\to A(K_j)$, as checked in
Theorem~\ref{B:kan-recovery}. Choose a lift of $Q$. Pulling it back
along $a_0$ chooses a lift of $P$; the trivial loop classes say that
pullback along every $a_j$ gives this same lift. Thus the whole
parallel family lifts to a family $\widetilde f$ in
$A(\widetilde K_j)$.

The finite complex $R_{\widetilde f}$ maps by its fat label into
$BA(\widetilde K_j)$, using (KC.2), and then into $BA_j$.
Its image in the weakly contractible $BA$ is nullhomotopic.
The finite-domain argument gives a further transition $j\to k$
for which this composite is already nullhomotopic in $BA_k$.
Apply the local claim to the functor
\[
 A(\widetilde K_j)\longrightarrow A(K_j)
       \longrightarrow A(K_k).
\]
Its source Kan complex is simply connected, its target is connected,
and its base map kills the indicated finite label. The image of
$\widetilde f$, which is the image of the original family, has an
integral affine equalizer. It remains an affine equalizer in $A$.

Every $A_i$ has common successors, by connectedness and the Kan
amalgamation proof in Theorem~\ref{B:kan-recovery}; this property
passes to the filtered colimit. We have proved the integral affine
equalizer criterion for every finite parallel family, so $\colim_A$
is integrally exact. The same integer witnesses give exactness over
every unital ring, including the zero ring. Integral exactness
implies totality by Theorem~\ref{F:companion-chain}, and totality
implies $BA\simeq *$ by applying its definition to the constant
singleton action. This proves all the equivalences.
\end{proof}
